
\def\PUDcodigo{ACE011}

\if\COMPILAR0
    \def\PUDcodacad{04507.13}
\else
    \def\PUDcodacad{04506.13}
\fi

\def\PUDdisciplina{Física 2}

\def\PUDsemestre{S3}

\def\PUDhoras{80}

%NOVOS ---------------------------------------------------
\def\PUDhorasTeoria{80}
\def\PUDhorasPratica{0}

\def\PUDinterECA{---}

\def\PUDatuaisECA{---}

\def\PUDrecursos{Projetor multimídia; Quadro branco e pincel.}

%----------------------------------------------------------

\def\PUDcreditos{4}

\def\PUDprerequisito{ \hyperref[PUD:ACE008]{ACE008} }

\if\COMPILAR0
    \def\PUDpreacad{ \hyperref[PUD:ACE008]{Física 1 (04507.8)} }
\else
    \def\PUDpreacad{ \hyperref[PUD:ACE008]{Física 1 (04506.8)} }
\fi

\def\PUDobjetivos{Analisar diversas situações físicas envolvendo rotação de corpos rígidos e a conservação do momento angular neste movimento, evidenciando suas aplicações para a engenharia.  Compreender o conceito das grandezas físicas envolvidas na estática e dinâmica dos fluidos, analisando problemas de hidrostática e hidrodinâmica baseados nas leis e princípios físicos envolvidos com aplicações para a engenharia. Compreender o movimento oscilatório, movimento harmônico simples, sua aplicação em pêndulos, o movimento harmônico simples amortecido, oscilações forçadas e ressonância inseridos em problemas práticos da engenharia. Conhecer os diversos tipos de ondas e as grandezas envolvidas na sua descrição, a equação de onda e o princípio da superposição para as ondas. Familiarizar-se com os principais conceitos relacionados à termodinâmica, às suas leis e aplicações na engenharia.}

\def\PUDementa{Cinemática e dinâmica da rotação;  Estática e dinâmica dos fluidos;  Oscilações;  Ondas;  Termodinâmica.}

\def\PUDconteudo{ &  Cinemática e Dinâmica da Rotação: Variáveis da rotação.  Relação entre as variáveis angulares e lineares.  Rotação com aceleração angular constante. 
Momento de Inércia.  Torque e segunda lei de Newton para rotação.  Trabalho e energia cinética da rotação.  Rolamento.  Momento angular.  Conservação do momento angular. 
Atividade prática em laboratório.  \\ 
 &  Estática e dinâmica dos fluidos: Massa específica e pressão.  Fluidos em repouso: teorema de Stevin, princípio de Pascal e princípio de Arquimedes. 
Fluidos ideais em movimento: equação da continuidade e equação de Bernoulli.  Atividade prática em laboratório.  \\ 
 &  Oscilações: Movimento harmônico simples.  Equações do MHS. 
Energia do MHS.  Pêndulos.  Movimento harmônico simples amortecido.  Oscilações forçadas e ressonância.  Atividade prática em laboratório.  \\ 
 &  Ondas: Tipos de ondas.  Ondas transversais e longitudinais.  Comprimento de onda e frequência.  Velocidade de uma onda.  Energia e potência de uma onda progressiva.  Equação de onda.  Princípio da superposição de ondas: interferência de ondas, ondas estacionárias e ressonância.  \\ 
 &  Termodinâmica: Temperatura, a lei zero da Termodinâmica e termômetros.  Dilatação térmica.  Temperatura e calor.  Primeira lei da termodinâmica. 
Mecanismos de transferência de calor e aplicações.  Processos irreversíveis e entropia.  Segunda lei da termodinâmica.  Aplicações a máquinas térmicas. }

\def\PUDmetodo{Aulas expositórias que podem ser teóricas e/ou práticas, onde as práticas no laboratório serão marcadas no decorrer da disciplina.}

\def\PUDavalia{A avaliação será feita com aplicação de provas teóricas e/ou práticas, além da possibilidade de inclusão de trabalhos e seminários no decorrer da disciplina.}

\def\PUDbibbasica{ &  \href{www.biblioteca.ifce.edu.br/index.asp?codigo_sophia=24544}{Young}, Hugh D.; Freedman, Roger A.  Física 2 - termodinâmica e ondas.  Volume 2.  12º ed.  São Paulo, SP: Addison-Wesley, 2008.  329p.  ISBN: 9788588639331. 
\\ 
 &   \href{www.biblioteca.ifce.edu.br/index.asp?codigo_sophia=65080}{Nussenzveig}, H. M.  Curso de Física 2 - Fluidos, Oscilações e Ondas. 5º ed.  São Paulo, SP: Edgard Blücher, 2014.  375p.  ISBN: 9788521207474. 
\\ 
 &   \href{www.biblioteca.ifce.edu.br/index.asp?codigo_sophia=63343}{Tipler}, Paul;  Mosca, Gene.  Fisica para cientistas e Engenheiros.  Volume 2.  6º ed.  Rio de Janeiro, RJ: LTC, 2015.  530p.  ISBN: 9788521617112. }

\def\PUDbibcomplementar{ &   \href{www.biblioteca.ifce.edu.br/index.asp?codigo_sophia=41355}{Halliday} e Resnick.  Fundamentos de Física 2 - gravitação, ondas e termodinâmica.  9º ed.  Rio de Janeiro, RJ: LTC, 2013.  296p.  ISBN: 9788521619048. 
\\ 
 &   \href{www.biblioteca.ifce.edu.br/index.asp?codigo_sophia=64975}{Junior}, F. R.; Nicolau, G. F.; Soares, P. A. T.  Os fundamentos da Física 2 - Termologia, Óptica e Ondas.  9º ed.  São Paulo, SP: Moderna, 2007.  532p.  ISBN: 9788516056575. 
\\ 
 &   \href{www.biblioteca.ifce.edu.br/index.asp?codigo_sophia=65410}{Borgnakke}, Claus.  Fundamentos da termodinâmica.  São Paulo, SP: Edgard Blücher, 2013.  728p.  ISBN: 9788521207924.
\\
 &  \href{www.biblioteca.ifce.edu.br/index.asp?codigo_sophia=8643}{Paraná}, Djalma Nunes.  Física.  E-book.  v.2.  São Paulo, SP: Ática, 1998.  ISBN: 85-08-070829.
\\
 &  \href{www.biblioteca.ifce.edu.br/index.asp?codigo_sophia=58221}{Sguazzardi}, Monica Midori Marcon Uchida.  Física Geral.  E-book.  Pearson.  140p.  ISBN: 9788543011080. }

\def\PUDautor{Samuel Vieira Dias}

\def\PUDdata{2014-07-28}

\def\PUDrevisao{3}

\def\PUDrevisor{Samuel Vieira Dias}

\def\PUDdatarevisao{2017-05-11}

\PUDcorpo

